% ===========================================================================
% Axioms for a Falling Universe: A Unified Field Theory of the RSVP Plenum
% Flyxion — Independent Researcher
% LuaLaTeX — Latin Modern Math
% ===========================================================================
\documentclass[12pt, a4paper]{article}

\usepackage{microtype}
\usepackage{geometry}
\geometry{a4paper, top=2.5cm, bottom=2.5cm, left=3cm, right=3cm}

\usepackage{amsmath, amsthm, amssymb}
\usepackage{mathtools}
\usepackage{physics}
\usepackage{hyperref}
\hypersetup{
  colorlinks=true,
  linkcolor=black,
  citecolor=black,
  urlcolor=black,
  pdftitle={Axioms for a Falling Universe},
  pdfauthor={Flyxion},
}
\usepackage{setspace}
\onehalfspacing

\usepackage{titlesec}
\titleformat{\section}{\large\bfseries}{\thesection.}{0.75em}{}
\titleformat{\subsection}{\normalsize\bfseries}{\thesubsection.}{0.75em}{}
\titleformat{\subsubsection}{\normalsize\itshape}{\thesubsubsection.}{0.75em}{}

\usepackage{abstract}
\renewcommand{\abstractnamefont}{\normalfont\bfseries}

\usepackage{thmtools}
\theoremstyle{plain}
\newtheorem{theorem}{Theorem}[section]
\newtheorem{proposition}[theorem]{Proposition}
\newtheorem{lemma}[theorem]{Lemma}
\newtheorem{corollary}[theorem]{Corollary}
\theoremstyle{definition}
\newtheorem{definition}[theorem]{Definition}
\newtheorem{axiom}[theorem]{Axiom}
\newtheorem{assumption}[theorem]{Assumption}
\theoremstyle{remark}
\newtheorem{remark}[theorem]{Remark}
\newtheorem{conjecture}[theorem]{Conjecture}

\usepackage{enumitem}
\setlist[enumerate]{nosep}

\usepackage{booktabs}
\usepackage{array}

% Notational shorthands
\newcommand{\Phif}{\Phi}
\newcommand{\vfield}{\mathbf{v}}
\newcommand{\Sfield}{S}
\newcommand{\LRSVP}{\mathcal{L}_{\mathrm{RSVP}}}
\newcommand{\ARSVP}{\mathcal{A}_{\mathrm{RSVP}}}
\newcommand{\Heff}{H_{\mathrm{eff}}}
\newcommand{\geff}{g_{\mu\nu}^{\mathrm{eff}}}
\newcommand{\Proj}[1]{\mathcal{P}_{#1}}
\newcommand{\Fij}{F_{ij}}
\newcommand{\divv}{\nabla\!\cdot\!\vfield}
\newcommand{\curlop}[1]{\nabla\!\times\!#1}
\newcommand{\Hspace}[1]{H^{#1}(\Omega)}
\newcommand{\Ltwo}{L^{2}(\Omega)}
\newcommand{\Linf}{L^{\infty}(\Omega)}
\newcommand{\BS}{B_S}
\newcommand{\lamsym}{\lambda_{\mathrm{sym}}}

\usepackage[numbers,sort&compress]{natbib}
\bibliographystyle{plainnat}

% ---------------------------------------------------------------------------
\begin{document}

\begin{titlepage}
\centering
\vspace*{2cm}
{\LARGE\bfseries Axioms for a Falling Universe:\par}
\vspace{0.6em}
{\Large A Unified Field Theory of the Relativistic Scalar--Vector Plenum\par}
\vspace{0.4em}
{\large From Lagrangian Uniqueness to Cognitive Closure\par}
\vspace{2.5cm}
{\large Flyxion\par}
{\normalsize Independent Researcher\par}
\vspace{0.5em}
{\normalsize May 2026\par}
\vspace{4cm}
\begin{abstract}
\noindent
We present a comprehensive and unified treatment of the Relativistic Scalar--Vector
Plenum (RSVP), an effective field theory of a continuous medium whose three coupled
fields --- the scalar potential $\Phi$, the lamphrodyne flow velocity $\vfield$, and the
configurational entropy density $S$ --- generate gravitational binding, cosmic
redshift, cognitive dynamics, semantic compression, and institutional closure as
multiscale coarse-grainings of a single variational principle.

The paper proceeds in four integrated movements.  First, we derive the unique
RSVP Lagrangian from seven kinematic and dynamical axioms by the logic of
effective field theory for continuous media, demonstrating that the divergence
constraint and the coupling structure are fixed up to higher-order corrections by
locality, rotational invariance, stability, and the constitutive identification of
gravity with entropic descent.  Second, we establish the complete PDE hierarchy
--- the RSVP Master System --- and prove local well-posedness, global weak
existence, Lyapunov stability, and exponential gradient decay using
semigroup methods and Aubin--Lions compactness.  Third, we show how a
projection ladder $\Proj{\mathrm{cosmo}} \to \Proj{\mathrm{semantic}} \to
\Proj{\mathrm{observer}} \to \Proj{\mathrm{institution}}$ maps the plenum
dynamics into the derived stacks EBSSC, Spherepop, TARTAN, and CLIO, each
preserving entropy monotonicity and sparsity invariants as coarse-graining
functors.  Fourth, we demonstrate that Jacobson's 1995 thermodynamic derivation
of the Einstein equation, Verlinde's entropic gravity, and Friston's free energy
principle all arise as special limits within this single architecture.  Falsifiable
predictions include an entropic redshift formula distinguishable from $\Lambda$CDM
at the $2\sigma$ level, coherence metrics for agency persistence in neuroimaging,
and institutional merge thresholds in historical data.

The theory is best characterised not as a cosmological alternative but as a
constrained nonequilibrium medium effective field theory whose geometric
coarse-graining induces effective spacetime curvature while entropy measures
admissible future trajectory volume.
\end{abstract}
\end{titlepage}

\newpage
\tableofcontents
\newpage

% ===========================================================================
\section{Introduction: The Problem of Fragmented Ontologies}
\label{sec:intro}
% ===========================================================================

Contemporary theoretical physics, cognitive science, and social theory share a structural
disability: each domain has developed its own ontological vocabulary, its own primitive
objects, and its own explanatory standards, and the cost of this disciplinary partition is
paid whenever a phenomenon crosses a boundary.  Cosmological expansion is described in
the language of Riemannian geometry; cognition in the language of neural computation;
institutional drift in the language of game theory.  When the same mathematical motif
--- gradient relaxation, entropy circulation, attractor formation --- appears in all three
domains, the standard response is to note the analogy with methodological caution and
proceed within the silo.

The present paper proposes a different response.  Rather than treating the recurrence of
the scalar--vector--entropy triplet $(\Phi, \vfield, S)$ across domains as an analogy to be
quarantined, we treat it as evidence of a shared substrate and attempt to make that
substrate mathematically precise.  The Relativistic Scalar--Vector Plenum (RSVP) is not
a metaphor for anything.  It is a candidate physical theory of a continuous medium whose
fields are geometric and thermodynamic structures, whose equations of motion are derived
from a variational principle, and whose coarse-graining functors produce the formal
objects that each disciplinary sub-theory already uses as primitives.

This is a stronger claim than unification by analogy, and it carries correspondingly stronger
obligations.  The paper discharges those obligations in four stages.  Section~\ref{sec:axioms}
derives the RSVP Lagrangian from seven axioms by the logic of effective field theory
for continuous media \cite{Nicolis2013, Endlich2013, Andersson2007}, proving that the
interaction structure is essentially unique at two-derivative order.  Section~\ref{sec:master}
establishes the RSVP Master System --- the coupled PDE hierarchy whose well-posedness we
prove using semigroup theory and energy methods.  Section~\ref{sec:stack} constructs the
projection ladder that connects the fundamental field equations to the derived stacks EBSSC,
Spherepop, TARTAN, and CLIO, formalising the sense in which these are distinct
coarse-grainings of the same plenum.  Section~\ref{sec:emergent} shows how Einstein gravity,
Verlinde's entropic force, and Friston's free energy principle all arise as limits.  Throughout,
falsifiable predictions are identified precisely so that the theory may be distinguished from
its competitors by observation.

Two preliminary clarifications are in order before proceeding.  First, the RSVP framework
is explicitly a candidate medium effective field theory, not a claim of final ontological
closure.  The seven axioms are physical selection principles, not metaphysical necessities.
Different axiom systems will produce different theories; the question is whether this
particular system is empirically viable.  Second, the paper distinguishes rigorously
between four distinct roles that equations play in the theory: kinematic axioms, which
constrain the geometry of the field space; dynamical constraints, which relate time
derivatives to spatial structure; constitutive relations, which identify physical observables
with field-theoretic quantities; and phenomenological identifications, which connect the
theory to measured quantities.  Conflating these roles by placing ``gravity equals entropic descent'' at the same logical level as locality obscures the distinction between kinematic constraints and constitutive interpretation. In the present formulation, that distinction is made explicit.

% ===========================================================================
\section{Fields, Symmetries, and the Axiom System}
\label{sec:axioms}
% ===========================================================================

\subsection{The Field Content}

The RSVP medium is described by three fields on a preferred $3+1$ foliation compatible
with medium-based effective field theories \cite{Andersson2007} and hydrodynamic gravity
analogues \cite{Barcelo2005}:

\begin{definition}[RSVP field triplet]
Let $\Omega \subset \mathbb{R}^3$ be an open bounded domain with Lipschitz boundary.
The RSVP fields are:
\begin{align}
\Phi &\in H^1(\Omega), && \text{scalar plenum potential (mass/meaning density),}\\
\vfield &\in [H^1(\Omega)]^3, && \text{lamphrodyne flow velocity,}\\
S &\in L^2(\Omega) \cap L^\infty(\Omega),\quad S\geq 0, && \text{configurational entropy density.}
\end{align}
\end{definition}

The scaling dimensions are assigned as $[\Phi] = \Delta_\Phi$, $[S] = \Delta_S$, and
$[\vfield_i] = \Delta_v$, with $[x] = -1$ in natural units.  The effective field theory
principle then determines which interaction terms are relevant, marginal, or irrelevant
under coarse-graining by power counting in $[\ell]^{-1}$ where $\ell$ is the infrared
cutoff scale \cite{Weinberg1995, Cheung2008}.

The distinction between ontic and epistemic interpretations of these fields is important
but should not be conflated with the derivational structure.  In the cosmological
projection, $\Phi$ functions as a physical potential, $\vfield$ as a bulk flow velocity,
and $S$ as a coarse-grained entropy per unit volume.  In the semantic projection, the
same fields represent meaning density, directed inference, and representational budget.
The projection ladder of Section~\ref{sec:stack} makes this reinterpretation map
precise.

\subsection{The Seven Axioms}

The selection of the RSVP Lagrangian from among all possible local field theories follows
the standard logic of effective field theory \cite{Weinberg1995}: one writes down the
most general action consistent with the symmetries and then truncates at the desired order
in derivatives.  The seven axioms below constitute the symmetry and stability requirements
that drive this selection.

\begin{axiom}[Locality and variational principle]
\label{ax:locality}
The dynamics derive from a local action functional
$\mathcal{A}[\Phi, \vfield, S] = \int \mathrm{d}t\int_\Omega \mathcal{L}(\Phi, \vfield, S, \partial_t\Phi, \nabla\Phi, \ldots)\,\mathrm{d}x$,
where $\mathcal{L}$ depends on fields and their first spacetime derivatives only.
Higher-derivative corrections are suppressed by the ultraviolet scale.
\end{axiom}

\begin{axiom}[Spatial rotational invariance]
\label{ax:rotation}
The Lagrangian density $\mathcal{L}$ is invariant under $SO(3)$ rotations of the
spatial arguments.  This forbids preferred-direction terms linear in $\vfield$ and
requires all vector contractions to appear as $|\vfield|^2$, $\nabla\Phi \cdot \vfield$,
$\nabla S \cdot \vfield$, or via the antisymmetric tensor $\epsilon_{ijk}$.
\end{axiom}

\begin{axiom}[Quadratic kinetic terms, ghost-freedom]
\label{ax:kinetic}
Each field possesses a canonically normalised kinetic term of the form
$\tfrac{1}{2}\dot{\Phi}^2$, $\tfrac{1}{2}|\dot{\vfield}|^2$, $\tfrac{1}{2}\dot{S}^2$.
Negative-norm kinetic terms (ghosts) are excluded by the requirement that the
classical phase space have a well-defined symplectic structure.
\end{axiom}

\begin{axiom}[Two-derivative truncation]
\label{ax:twoderiv}
At long wavelengths the EFT is truncated at two spatial derivatives per field.
All terms containing three or more spatial derivatives are absorbed into
higher-order corrections suppressed by the ratio of the infrared to ultraviolet
scales.
\end{axiom}

\begin{axiom}[Energy bounded below]
\label{ax:stability}
The Hamiltonian density $\mathcal{H}$ derived from $\mathcal{L}$ by Legendre
transformation satisfies $\mathcal{H} \geq 0$ pointwise for all field configurations.
This ensures Lyapunov stability of the ground state and excludes runaway solutions.
\end{axiom}

\begin{axiom}[Lamphrodyne divergence constraint]
\label{ax:constraint}
The medium is subject to the constraint
\begin{equation}
\divv \approx \alpha_\Phi \Phi + \alpha_S S,
\label{eq:constraint}
\end{equation}
analogous to incompressibility conditions in continuum EFTs \cite{Nicolis2013} but
modified by density and entropy sources.  This is a kinematic axiom: it restricts
the configuration space of admissible field profiles without specifying dynamics.
\end{axiom}

\begin{axiom}[Gravity as entropic descent --- constitutive relation]
\label{ax:gravity}
\emph{(Constitutive, not kinematic.)} Gravitational attraction is identified with the
entropic descent of flow trajectories: the force on a test particle is proportional
to the entropy gradient $\nabla S$ evaluated along the flow.  This axiom generalises
the Clausius--Raychaudhuri logic \cite{Jacobson1995} and the entropic force
programme \cite{Verlinde2011} and is imposed as a phenomenological closure condition
rather than a primitive assumption of the same logical rank as
Axioms~\ref{ax:locality}--\ref{ax:stability}.
\end{axiom}

Axiom~\ref{ax:gravity} differs from the preceding axioms in logical role. Whereas Axioms~\ref{ax:locality}--\ref{ax:constraint} determine the admissible structure of the Lagrangian, Axiom~\ref{ax:gravity} provides a constitutive interpretation of the resulting field dynamics. The Lagrangian itself is fully determined by Axioms~\ref{ax:locality}--\ref{ax:constraint}; Axiom~\ref{ax:gravity} is introduced subsequently to relate the entropic flow structure of the theory to emergent gravitational behaviour.

\subsection{Uniqueness of the Lagrangian}

\begin{theorem}[Uniqueness of the RSVP Lagrangian]
\label{thm:uniqueness}
Under Axioms~\ref{ax:locality}--\ref{ax:constraint}, the unique local, rotationally
invariant, ghost-free, two-derivative, energy-bounded Lagrangian for the fields
$(\Phi, \vfield, S)$ coupled via the constraint~\eqref{eq:constraint} is, up to
boundary terms and higher-order corrections,
\begin{align}
\LRSVP &= \tfrac{1}{2}\dot{\Phi}^2 - \tfrac{c_\Phi^2}{2}|\nabla\Phi|^2 - U_\Phi(\Phi)
\notag\\
&\quad + \tfrac{1}{2}|\dot{\vfield}|^2 - \tfrac{c_v^2}{4}F_{ij}F^{ij}
       - \tfrac{\kappa_v}{2}\!\left(\divv - \alpha_\Phi\Phi - \alpha_S S\right)^2
\notag\\
&\quad + \tfrac{1}{2}\dot{S}^2 - \tfrac{c_S^2}{2}|\nabla S|^2 - U_S(S)
\notag\\
&\quad + g_1\,\Phi\,\vfield\cdot\nabla S
       + g_2\,S\,\vfield\cdot\nabla\Phi
       + g_3\,\Phi S,
\label{eq:Lagrangian}
\end{align}
where $F_{ij} = \partial_i v_j - \partial_j v_i$ and $c_\Phi, c_v, c_S, \kappa_v, \alpha_\Phi, \alpha_S, g_1, g_2, g_3$ are
positive coupling constants.
\end{theorem}

\begin{proof}
We enumerate the allowed operators at each order in derivatives, compatible with $SO(3)$
invariance and the field content $(\Phi, \vfield, S)$.

\emph{Zeroth order in derivatives (potential terms).}  The most general $SO(3)$-invariant
potential is an arbitrary smooth function $V(\Phi, S, |\vfield|^2)$.  The ghost-freedom
axiom requires this to be bounded below; stability then forces $V$ to admit minima. We
split $V = U_\Phi(\Phi) + U_S(S) + V_{\mathrm{int}}$ where $V_{\mathrm{int}}$ contains
cross terms.  At quadratic order in the fields, $V_{\mathrm{int}} \supset g_3 \Phi S$.
Terms proportional to $|\vfield|^2$ at zeroth order in gradients are absorbed into the
kinetic sector by field redefinition.

\emph{First order in derivatives.}  No $SO(3)$-invariant scalar linear in $\nabla\Phi$,
$\nabla S$, or $\vfield$ can be constructed without a fixed preferred vector, which
Axiom~\ref{ax:rotation} excludes.  Cross terms of the form $\vfield\cdot\nabla\Phi$ and
$\vfield\cdot\nabla S$ are first-order in derivatives; they are permitted by symmetry.
The $SO(3)$-invariant combinations are $\Phi\,\vfield\cdot\nabla S$ and
$S\,\vfield\cdot\nabla\Phi$, giving the $g_1$ and $g_2$ couplings.

\emph{Two-derivative kinetic terms.}  The ghost-free, rotationally invariant kinetic terms
are $\tfrac{1}{2}(\partial_t X)^2 - \tfrac{c_X^2}{2}|\nabla X|^2$ for each scalar
$X \in \{\Phi, S\}$.  For the vector field, rotational invariance permits two independent
two-derivative structures: $|\partial_t \vfield|^2$ and the combination
$\partial_i v_j \partial^i v^j - \partial_i v_j \partial^j v^i = \tfrac{1}{2}F_{ij}F^{ij}$ (the
curl-squared).  The longitudinal two-derivative structure $(\divv)^2$ is penalised by
the constraint~\eqref{eq:constraint}, which enters as the Lagrange multiplier term
$-\tfrac{\kappa_v}{2}(\divv - \alpha_\Phi\Phi - \alpha_S S)^2$. In the stiff limit
$\kappa_v \to \infty$ this enforces the constraint exactly.

\emph{Completeness.}  No further $SO(3)$-invariant, two-derivative, ghost-free operators
can be formed from $(\Phi, \vfield, S)$ without violating the truncation axiom.  The
Lagrangian~\eqref{eq:Lagrangian} is therefore unique up to the values of the coupling
constants.
\end{proof}

The divergence penalty term $-\tfrac{\kappa_v}{2}(\divv - \alpha_\Phi\Phi - \alpha_S S)^2$
deserves special attention because it performs multiple conceptual functions simultaneously.
It acts as a medium-compressibility condition, an emergent curvature source, a geometric
admissibility relation, and a coupling between the longitudinal vector mode and the scalar
and entropy fields.  In the stiff limit it eliminates the longitudinal mode entirely,
projecting onto transverse (divergence-free) flow.  Many consequences of the theory
--- the Jacobson derivation, the Verlinde limit, the CLIO descent algorithm --- ultimately
trace back to this single term.

% ===========================================================================
\section{The RSVP Master System: PDE Formulation and Well-Posedness}
\label{sec:master}
% ===========================================================================

\subsection{Derivation of the Field Equations}

\begin{theorem}[Euler--Lagrange equations of RSVP]
\label{thm:EL}
Stationary variation of the action $\ARSVP = \int \mathrm{d}t\int_\Omega \LRSVP\,\mathrm{d}x$
with homogeneous Neumann boundary conditions $\partial_n\Phi = \partial_n S = \partial_n\vfield = 0$
on $\partial\Omega$ yields the RSVP field equations:
\begin{align}
\ddot{\Phi} - c_\Phi^2\Delta\Phi + U'_\Phi(\Phi)
  &= \kappa_v\alpha_\Phi(\divv - \alpha_\Phi\Phi - \alpha_S S)
    + g_1\vfield\cdot\nabla S + g_2 S(\divv) + g_3 S,
\label{eq:EL-Phi}\\
\ddot{\vfield} + c_v^2\nabla\!\times\!(\nabla\!\times\!\vfield)
  &= \kappa_v(\divv - \alpha_\Phi\Phi - \alpha_S S)\nabla(\divv)
    + g_1\Phi\nabla S + g_2 S\nabla\Phi,
\label{eq:EL-v}\\
\ddot{S} - c_S^2\Delta S + U'_S(S)
  &= \kappa_v\alpha_S(\divv - \alpha_\Phi\Phi - \alpha_S S)
    + g_2\vfield\cdot\nabla\Phi + g_1\nabla\cdot(\Phi\vfield) + g_3\Phi.
\label{eq:EL-S}
\end{align}
\end{theorem}

\begin{proof}
The variation $\delta_\Phi\ARSVP = 0$ gives
\[
\ddot{\Phi} - c_\Phi^2\Delta\Phi + U'_\Phi(\Phi)
+ \kappa_v\alpha_\Phi(\alpha_\Phi\Phi + \alpha_S S - \divv)
- g_1\vfield\cdot\nabla S - g_2 S\nabla\cdot\vfield - g_3 S = 0.
\]
Rearranging and defining $\chi \equiv \divv - \alpha_\Phi\Phi - \alpha_S S$ yields
\eqref{eq:EL-Phi}.  The variation $\delta_v\ARSVP = 0$ expands the curl term via
$\nabla\times(\nabla\times\vfield) = \nabla(\divv) - \Delta\vfield$; using $\kappa_v\chi\nabla\chi$ one
obtains \eqref{eq:EL-v}.  The variation $\delta_S\ARSVP = 0$ proceeds analogously,
noting that $g_1\nabla\cdot(\Phi\vfield) = g_1\vfield\cdot\nabla\Phi + g_1\Phi(\divv)$
after integration by parts.
\end{proof}

\subsection{The Overdamped Master System}

The second-order system \eqref{eq:EL-Phi}--\eqref{eq:EL-S} becomes the physically most
tractable form in the overdamped (dissipative) limit, where inertial terms $\ddot{X}$ are
negligible relative to diffusion and damping.  This limit is natural for the semantic and
cognitive projections where irreversibility dominates.  Including the Rayleigh dissipation
functional and passing to the overdamped limit produces the RSVP Master System:

\begin{definition}[RSVP Master System]
\label{def:master}
The RSVP Master System is the coupled PDE system on $\Omega\times(0,\infty)$:
\begin{align}
\partial_t\Phi &= -\nabla\cdot(\Phi\vfield) + \kappa_\Phi\Delta\Phi + \sigma S - U'_\Phi(\Phi),
\label{eq:master-Phi}\\
\partial_t\vfield &= -(\vfield\cdot\nabla)\vfield - \lambda\nabla\Phi - \nu\vfield
                    + \kappa_v(\nabla\times\vfield) + \eta,
\label{eq:master-v}\\
\partial_t S &= D_S\Delta S - \mu\Phi + \chi|\vfield|^2 + \xi(t),
\label{eq:master-S}
\end{align}
with initial data $(\Phi_0, \vfield_0, S_0)$ satisfying Definition~1, homogeneous Neumann
boundary conditions, and parameters $\kappa_\Phi, D_S, \lambda, \nu, \mu, \sigma, \chi > 0$.
Here $\eta$ is a stochastic forcing term and $\xi(t)$ a thermal noise process, both
assumed to satisfy standard Gaussian conditions with bounded variance.
\end{definition}

Each term in this system has a precise physical interpretation.  In
\eqref{eq:master-Phi}: the divergence term $-\nabla\cdot(\Phi\vfield)$ is advection of
scalar density by the flow; $\kappa_\Phi\Delta\Phi$ is diffusive smoothing (lamphrodyne
relaxation); $\sigma S$ injects negentropy into the potential from the entropy reservoir;
and $-U'_\Phi(\Phi)$ provides attractor geometry.  In \eqref{eq:master-v}: the
nonlinear term $-(\vfield\cdot\nabla)\vfield$ is the Euler convective derivative;
$-\lambda\nabla\Phi$ drives flow down the potential gradient; $-\nu\vfield$ is viscous
damping; $\kappa_v(\nabla\times\vfield)$ introduces vorticity-driven circulation; and
$\eta$ models stochastic excitation.  In \eqref{eq:master-S}: $D_S\Delta S$ diffuses
entropy; $-\mu\Phi$ expresses that high potential suppresses local entropy (negentropy
concentration); $\chi|\vfield|^2$ sources entropy from kinetic activity; and $\xi(t)$
is thermal noise.

\subsection{Well-Posedness}

\begin{assumption}
\label{ass:regularity}
The potential $U_\Phi \in C^2(\mathbb{R})$ satisfies $U''_\Phi \geq 0$.  The source functions
satisfy $\sigma(\Phi,\vfield) \leq C(1 + |\Phi| + |\vfield|^2)$ for some constant $C > 0$.
The diffusion coefficients satisfy $\kappa_\Phi, D_S, \kappa_v > 0$.  Noise terms $\eta$
and $\xi$ are adapted, square-integrable, and of bounded variance.
\end{assumption}

\begin{theorem}[Local classical well-posedness]
\label{thm:local}
Under Assumption~\ref{ass:regularity}, there exists $T^* > 0$ depending on
$\|\Phi_0\|_{H^1}$, $\|\vfield_0\|_{H^1}$, and $\|S_0\|_{L^\infty}$ such that the RSVP
Master System admits a unique classical solution
\[
\Phi \in C([0,T^*]; H^2(\Omega)) \cap C^1([0,T^*]; H^1(\Omega)),
\]
and similarly for $\vfield$ and $S$.
\end{theorem}

\begin{proof}
Rewrite the system as $\partial_t u = Lu + \mathcal{N}(u)$ where $u = (\Phi, \vfield, S)$
and $L$ is the sectorial operator generated by $(-\kappa_\Phi\Delta, -\kappa_v\Delta, -D_S\Delta)$
with Neumann boundary conditions.  By standard spectral theory, $L$ generates an analytic
semigroup $\{e^{tL}\}_{t\geq 0}$ on $H^1(\Omega)\times [H^1(\Omega)]^3 \times L^2(\Omega)$.
The nonlinearity $\mathcal{N}$ is locally Lipschitz from $H^2\times[H^2]^3\times H^1$ into
$H^1\times[H^1]^3\times L^2$ by the Sobolev embedding $H^1(\Omega)\hookrightarrow L^6(\Omega)$
(in $d = 3$) and the polynomial growth condition in Assumption~\ref{ass:regularity}.
The abstract Cauchy theorem for analytic semigroups \cite{Pazy1983} then yields local existence
and uniqueness.
\end{proof}

\begin{theorem}[Global weak solutions]
\label{thm:global}
Under Assumption~\ref{ass:regularity}, there exists a global weak solution in the
Leray--Hopf sense:
\[
\Phi \in L^\infty(0,\infty; H^1) \cap L^2_{\mathrm{loc}}(0,\infty; H^2),
\quad
\partial_t\Phi \in L^2_{\mathrm{loc}}(0,\infty; H^{-1}),
\]
and analogously for $\vfield$ and $S$.
\end{theorem}

\begin{proof}
Construct a Galerkin approximation using eigenfunctions $\{e_k\}$ of $-\Delta$ with
Neumann boundary conditions.  The energy estimate
\[
\frac{\mathrm{d}}{\mathrm{d}t}\left(\frac{1}{2}\int_\Omega|\Phi|^2
+ \frac{1}{2}\int_\Omega|\vfield|^2 + \int_\Omega S\right) \leq C,
\]
follows from the growth condition on $\sigma$ and Young's inequality.  Uniform bounds on
the Galerkin approximants in the indicated spaces follow from the Poincaré inequality.
Strong convergence in $L^2(0,T;L^2)$ is obtained via the Aubin--Lions compactness lemma
\cite{Lions1969}, which provides the compactness needed to pass to the limit in the
nonlinear terms.
\end{proof}

\subsection{Entropy Bounds and Lyapunov Structure}

\begin{lemma}[Global entropy bound]
\label{lem:entropybound}
Under Assumption~\ref{ass:regularity}, if $\sigma(\Phi,\vfield) \leq K(|\Phi| + |\vfield|^2)$
and $\mu > 0$, then
\[
\int_\Omega S(x,t)\,\mathrm{d}x \leq B_S
:= \max\!\left\{\int_\Omega S_0,\;\frac{\mathrm{ess\,sup}\;\sigma}{\mu}|\Omega|\right\}.
\]
\end{lemma}

\begin{proof}
Integrate \eqref{eq:master-S} over $\Omega$ and apply the Neumann boundary condition:
\[
\frac{\mathrm{d}}{\mathrm{d}t}\int_\Omega S = \int_\Omega \sigma(\Phi,\vfield)
- \mu\int_\Omega\Phi + \chi\int_\Omega|\vfield|^2.
\]
The term $-\mu\int\Phi$ controls entropy decay; estimating $\int\sigma \leq K(\int|\Phi| + \int|\vfield|^2)$
and applying Grönwall's inequality yields the stated bound.
\end{proof}

\begin{theorem}[Lyapunov functional]
\label{thm:lyapunov}
Define the entropy Lyapunov functional $\mathcal{H}(t) = \int_\Omega S(x,t)\,\mathrm{d}x$.
Then $\dot{\mathcal{H}} \leq \mathrm{ess\,sup}\;\sigma - \mu\mathcal{H}$, so
$\mathcal{H}(t)$ is uniformly bounded and approaches $B_S/\mu$ asymptotically.
\end{theorem}

\begin{theorem}[Exponential gradient decay in the lamphrodyne regime]
\label{thm:gradientdecay}
If $\sigma \equiv 0$ on $[T,\infty)$, then
\[
\|\nabla S(t)\|_{L^2} \leq \|\nabla S(T)\|_{L^2}\,e^{-\mu_0(t-T)},
\qquad \mu_0 = \min(\mu, D_S\lambda_1),
\]
where $\lambda_1 > 0$ is the first nonzero Neumann eigenvalue of $-\Delta$.
\end{theorem}

\begin{proof}
Multiply \eqref{eq:master-S} (with $\sigma = 0$) by $-\Delta S$ and integrate over $\Omega$:
\[
\frac{1}{2}\frac{\mathrm{d}}{\mathrm{d}t}\int_\Omega|\nabla S|^2
+ D_S\int_\Omega|\Delta S|^2 + \mu\int_\Omega|\nabla S|^2 = 0.
\]
The Poincaré inequality gives $\int|\Delta S|^2 \geq \lambda_1\int|\nabla S|^2$, from
which $\tfrac{1}{2}\dot{y} + (D_S\lambda_1 + \mu)y \leq 0$ where
$y = \|\nabla S\|_{L^2}^2$, yielding the exponential bound.
\end{proof}

This theorem formalises the ``lamphrodyne smoothing'' property: the entropy field
exponentially relaxes toward spatial uniformity in the absence of sources.  The timescale
$\mu_0^{-1}$ sets the characteristic coarse-graining time of the medium.

\subsection{Hamiltonian and Dirac-Bracket Formulation in the Stiff Limit}

In the stiff limit $\kappa_v \to \infty$ the divergence constraint \eqref{eq:constraint}
becomes exact: $\chi \equiv \divv - \alpha_\Phi\Phi - \alpha_S S = 0$.  This converts
the constraint from a soft penalty into a genuine second-class constraint in the
Dirac--Bergmann sense \cite{Dirac1964, HenneauxTeitelboim1992}.

\begin{theorem}[Dirac bracket in the stiff limit]
\label{thm:dirac}
In the stiff limit, the canonical momenta satisfy the constraint $\chi = 0$ as a
second-class constraint.  The Dirac bracket between the velocity field and its conjugate
momentum is
\[
\{v_i(\mathbf{x}), \pi^j_v(\mathbf{y})\}^* = P^j_i(\mathbf{x} - \mathbf{y}),
\]
where $P^j_i(\mathbf{z}) = \delta^j_i\delta^{(3)}(\mathbf{z})
- \partial_i\partial^j G(\mathbf{z})$ and $G$ is the Neumann Green's function of
$-\Delta$ on $\Omega$.  This projects onto transverse (divergence-free) flow modes.
\end{theorem}

\begin{proof}
Following Dirac's algorithm \cite{Dirac1964}, the primary constraint is $\chi = 0$ and
the secondary constraint is $\{\chi, H\}_{\mathrm{PB}} = 0$, which is satisfied automatically
by the form of $H$.  The constraint matrix $C_{ab} = \{\chi_a, \chi_b\}$ is invertible,
and the Dirac bracket is $\{F,G\}^* = \{F,G\}_{\mathrm{PB}} - \{F,\chi_a\}(C^{-1})^{ab}\{\chi_b,G\}$.
The kernel of this bracket on velocity-momentum pairs gives the transverse projector
$P^j_i$ via the formula for the Green's function of the Laplacian with Neumann conditions.
\end{proof}

The transverse projector in Theorem~\ref{thm:dirac} has a precise physical interpretation:
it is the operator that removes all compression from the flow, leaving only the
divergence-free, ``vortical'' degrees of freedom.  This is the medium-theoretic analogue
of the Coulomb gauge in electrodynamics.  The stiff-limit RSVP theory is therefore
a theory of incompressible, entropy-modulated vortex flow --- a constrained Hamiltonian
system with a well-developed formal machinery \cite{Nicolis2013, Endlich2013}.

% ===========================================================================
\section{The Projection Ladder: From Plenum to Derived Stacks}
\label{sec:stack}
% ===========================================================================

The central claim of RSVP is that cosmology, cognition, and institutional dynamics arise as distinct coarse-grainings of the same underlying plenum. In this view, the differences between physical, semantic, and social systems do not originate from fundamentally different substrates, but from different projection scales, admissibility constraints, and informational resolutions imposed upon a shared field structure. Observable structures at each level therefore emerge not as isolated ontologies, but as effective descriptions induced by successive reductions of the same scalar--vector--entropy dynamics.

This section constructs the projection ladder connecting these levels of description and formalises the corresponding coarse-graining operators. It then establishes that key invariants --- particularly entropy monotonicity and sparsity structure --- are preserved under projection, ensuring that the derived effective systems remain dynamically consistent with the underlying plenum evolution.

\subsection{The Projection Ladder}

\begin{definition}[Projection operators]
\label{def:projections}
Define the following coarse-graining projections of the RSVP field triplet
$(\Phi, \vfield, S)$:
\begin{align}
\Proj{\mathrm{cosmo}} &: (\Phi, \vfield, S) \mapsto (\rho_{\mathrm{eff}}, g_{\mu\nu}^{\mathrm{eff}}, \mathcal{T})
&&\text{(gravitational/cosmological observables),}\\
\Proj{\mathrm{semantic}} &: (\Phi, \vfield, S) \mapsto (\mathcal{M}, \pi^*, H_{\mathrm{sem}})
&&\text{(EBSSC sparse semantic manifold),}\\
\Proj{\mathrm{observer}} &: (\Phi, \vfield, S) \mapsto (\mathcal{O}_n, \Gamma, \mathcal{C})
&&\text{(CLIO holographic observer),}\\
\Proj{\mathrm{institution}} &: (\Phi, \vfield, S) \mapsto (\mathcal{S}_i, \mathcal{R}, \Sigma)
&&\text{(Spherepop institutional sphere).}
\end{align}
Each $\mathcal{P}$ is a bounded linear map on the relevant Sobolev space that commutes
with the time evolution generated by the Master System.
\end{definition}

The key claim of Definition~\ref{def:projections} is that each projection commutes with
the RSVP dynamics.  This is not automatic; it requires that the coarse-graining map
$\mathcal{P}$ satisfies $\mathcal{P}(\partial_t X) = \partial_t(\mathcal{P}(X))$,
which holds when $\mathcal{P}$ is given by convolution with a smoothing kernel $K$ or
conditional expectation onto a coarser sigma-algebra.

\begin{theorem}[Entropy monotonicity under coarse-graining]
\label{thm:entropymonotone}
Let $\mathcal{P}$ be any coarse-graining projection defined by conditional expectation
onto a coarser sigma-algebra $\mathcal{G}$.  Then
\[
H(\mathcal{P}(S)) \leq H(S) + C_{\mathcal{P}},
\]
where $H$ denotes Shannon entropy, the inequality reflects the data processing
inequality, and $C_{\mathcal{P}}$ is a bounded constant depending only on the
coarsening level.
\end{theorem}

\begin{proof}
The coarse-graining is $\mathcal{P}(S) = \mathbb{E}[S\,|\,\mathcal{G}]$ for a coarser
sigma-algebra $\mathcal{G}$.  By Jensen's inequality applied to the convex function
$x\mapsto -x\log x$:
\[
H(\mathcal{P}(S)) = H(\mathbb{E}[S|\mathcal{G}]) \leq \mathbb{E}[H(S|\mathcal{G})] \leq H(S).
\]
The additive constant $C_{\mathcal{P}}$ accounts for the entropy introduced by the
coarsening noise, which is bounded by the log-cardinality of the coarsening partition.
\end{proof}

\subsection{EBSSC: Entropy-Bounded Sparse Semantic Calculus}

The semantic projection $\Proj{\mathrm{semantic}}$ maps the plenum onto the
entropy-constrained sparse representation manifold.

\begin{definition}[Semantic manifold]
\label{def:ebssc}
The Entropy-Bounded Sparse Semantic Calculus (EBSSC) is the triplet
$\mathcal{M} = (\Phi, \vfield, S)$ together with the sparsity prior
$\|\Phi\|_0 \ll N$ (where $N = |\Omega|$ in discretised form) and the optimal sparse
policy
\[
\pi^* = \arg\min_{\pi}\,\mathbb{E}\!\left[S + \alpha\|\Phi\|_0 + \beta\|\divv\|^2\right].
\]
\end{definition}

\begin{theorem}[Existence of optimal sparse policy]
\label{thm:sparse}
Under the assumption that the action space is finite and $\Phi \in \ell^0$, there
exists a minimiser $\pi^* \in \mathcal{P}(\mathcal{A})$ satisfying the EBSSC
optimisation problem.
\end{theorem}

\begin{proof}
The functional $J(\pi) = \mathbb{E}_\pi[S] + \alpha\|\Phi\|_0 + \beta\mathbb{E}_\pi[\|\divv\|^2]$
is lower semi-continuous on the compact simplex $\mathcal{P}(\mathcal{A})$ in the
weak-$*$ topology (by Portmanteau).  The $\ell^0$ term is lower semi-continuous on
finite-dimensional spaces because level sets $\{\|\Phi\|_0 \leq k\}$ are finite unions
of coordinate subspaces, hence closed.  Existence follows from the extreme value theorem.
\end{proof}

\begin{theorem}[Sparsity--entropy tradeoff]
\label{thm:tradeoff}
At the optimum $\pi^*$,
\[
\|\Phi\|_0 \leq \frac{\mathbb{E}[S] + \beta\,\mathbb{E}[\|\divv\|^2]}{\alpha}.
\]
\end{theorem}

\begin{proof}
At the optimum, $J(\pi^*) \geq \alpha\|\Phi\|_0$ since $\mathbb{E}[S]$ and
$\beta\mathbb{E}[\|\divv\|^2]$ are both non-negative.  Rearranging gives the bound.
\end{proof}

\subsection{Spherepop Calculus: Institutions as Bounded Thermodynamic Spheres}

The institutional projection $\Proj{\mathrm{institution}}$ maps the plenum onto a
lattice of thermodynamic spheres whose merger dynamics are controlled by entropy
thresholds and ritual resistance.

\begin{definition}[Institutional sphere]
An institutional sphere is a triple $\mathcal{S} = (\mathcal{I}, \mathcal{B}, \Sigma)$
where $\mathcal{I}$ is the interior field configuration, $\mathcal{B}$ is the boundary
operator, and $\Sigma \in [0,1]$ is the permeability coefficient.  The Pop operation
is $\mathcal{P}: \mathcal{S}_i \otimes \mathcal{S}_j \to \mathcal{S}_k$, and merge
occurs when $\Sigma_i + \Sigma_j > \Theta_{\mathrm{closure}}$.
\end{definition}

\begin{theorem}[Merge condition and resistance]
\label{thm:merge}
Merge succeeds if and only if $\Sigma_i + \Sigma_j > \Theta_{\mathrm{closure}}$ and
$R_t + 2H < P_{\max}$, where $R_t = \int_0^T\gamma(t)\,\mathrm{d}t$ is the accumulated
ritual cost and $H$ is the keyspace entropy.  The resistance grows at rate
$\dot{R} \geq \gamma_{\min} + \log_2(1 + \delta H(t))$ for incremental hardening.
\end{theorem}

\begin{theorem}[Non-flattening merge preserves volume]
For $\mathcal{S}_a \oplus \mathcal{S}_b$,
$|\mathcal{I}_a \times \mathcal{I}_b| = |\mathcal{I}_a|\cdot|\mathcal{I}_b|$
and $\Sigma_\oplus = \min(\Sigma_a,\Sigma_b)$.
\end{theorem}

\subsection{TARTAN: Recursive Memory Lattices}

\begin{definition}[TARTAN tile]
A TARTAN (Trajectory-Aware Recursive Tiling with Annotated Noise) memory tile is a
quadruple $\mathcal{T} = (\Phi_\mathcal{T}, \vfield_\mathcal{T}, S_\mathcal{T}, m_\mathcal{T})$
where $m_\mathcal{T}$ is an annotation vector.  The refinement operator is
$f(\mathcal{T}) = \mathrm{refine}(\mathcal{T}) + \mathrm{annotate}(\mathcal{T}) + \mathrm{smooth}(\mathcal{T})$.
Semantic noise is injected as $\eta_s \sim \mathcal{N}(0, c|\nabla\Phi|^2)$.
\end{definition}

\begin{theorem}[Fixed-point convergence of TARTAN refinement]
\label{thm:tartan}
If the refinement operator $f$ is contractive in $L^2$ with constant $\rho < 1$, then
the lattice sequence converges:
$\|\mathcal{T}_n - \mathcal{T}^*\|_{L^2} \leq \rho^n\|\mathcal{T}_0 - \mathcal{T}^*\|_{L^2}$.
\end{theorem}

\begin{proof}
Direct application of the Banach contraction mapping theorem on the complete metric space
$(L^2(\Omega), \|\cdot\|_{L^2})$.
\end{proof}

\begin{theorem}[Gradient-weighted noise maximises boundary exploration]
The variance $c|\nabla\Phi|^2$ of the semantic noise process maximises expected Fisher
information in a neighbourhood of decision boundaries, concentrating exploration where
the posterior has highest curvature.
\end{theorem}

\begin{proof}
By information geometry, the Fisher information at parameter $\theta$ is
$\mathcal{I}(\theta) = \mathbb{E}[\|\nabla_\theta\log p\|^2]$.  The noise variance
$\sigma_s^2(\nabla\Phi) = c|\nabla\Phi|^2$ is proportional to the squared gradient of the
log-density, so $\mathbb{E}[\mathcal{I}(\theta)\sigma_s^2] \propto \int|\nabla\Phi|^2|\nabla\log p|^2$,
which is maximised by aligning $\sigma_s$ with $|\nabla\Phi|$.
\end{proof}

\subsection{CLIO: Cognitive Descent as Plenum Projection}

\begin{definition}[CLIO descent]
The Constrained Learning and Inference Optimiser (CLIO) is the gradient flow
\[
\Phi_{t+1} = \Phi_t - \eta_t\nabla_\Phi\mathcal{L}(\Phi_t, \vfield_t, S_t),
\qquad
\mathcal{L} = \mathbb{E}[S] + \lambda_{\mathrm{sparse}}\|\Phi\|_0 - \lambda_{\mathrm{flow}}\|\vfield\|^2.
\]
\end{definition}

The CLIO loss functional is a direct projection of the RSVP action: the entropy term
penalises disorder, the sparsity term rewards compressed representations, and the flow
reward $-\lambda_{\mathrm{flow}}\|\vfield\|^2$ prevents collapse to zero agency.

\begin{theorem}[Linear convergence of CLIO descent]
\label{thm:clio}
If $\mathcal{L}$ is $\mu$-strongly convex in $\Phi$ on sparse supports and
$\nabla_\Phi\mathcal{L}$ is $L$-Lipschitz, then with $\eta_t = 1/L$:
\[
\mathcal{L}(\Phi_t) - \mathcal{L}^* \leq \left(1 - \frac{\mu}{L}\right)^t(\mathcal{L}(\Phi_0) - \mathcal{L}^*).
\]
\end{theorem}

\begin{theorem}[Flow reward prevents agency collapse]
\label{thm:agency}
If $\lambda_{\mathrm{flow}} > 0$, then $\vfield_t \not\to 0$ along CLIO trajectories
unless $\Phi$ is identically flat.
\end{theorem}

\begin{proof}
Suppose $\vfield_t \to 0$.  Then $\partial_t\mathcal{L} \leq \mathbb{E}[S] + \lambda_{\mathrm{sparse}}\|\Phi\|_0$.
But from \eqref{eq:master-v}, $\vfield$ is driven by $-\lambda\nabla\Phi$, so $\vfield\to 0$
requires $\nabla\Phi\to 0$, i.e.\ $\Phi$ is asymptotically constant.  The flow reward
then contributes $-\lambda_{\mathrm{flow}}\|\vfield\|^2 \to 0$, but the initial contribution
$-\lambda_{\mathrm{flow}}\|\vfield_0\|^2 < 0$ means the loss is strictly lower than the
$\vfield = 0$ level, contradicting minimality unless $\Phi$ is flat.
\end{proof}

Theorem~\ref{thm:agency} captures a deep feature of the RSVP architecture: the flow field
$\vfield$ is not a passive tracer of the scalar gradient but an active component of the
system with its own dynamics.  Setting $\lambda_{\mathrm{flow}} = 0$ decouples this
agency and allows the system to collapse to entropic equilibrium with no directed action.

\subsection{Agency Coherence and Observer Holography}

The observer projection $\Proj{\mathrm{observer}}$ maps the plenum onto a sequence of
holographic images through weighted inner products:
\[
\mathcal{O}_n(\Phi) = \int_\Omega \Phi(\mathbf{x})\,w_n(\mathbf{x})\,\mathrm{d}\mathbf{x},
\]
where $w_n$ is the $n$-th perceptual kernel.

\begin{definition}[Agency coherence metric]
The agency coherence metric is
\[
\Gamma(t) = \frac{I(\Phi;\vfield)}{H(S)},
\]
where $I(\Phi;\vfield)$ is the mutual information between scalar and vector fields in
local patches, and $H(S)$ is the Shannon entropy of the entropy field distribution.
\end{definition}

\begin{theorem}[Upper bound on coherence]
Under the joint Gaussian assumption on field fluctuations,
\[
\Gamma(t) \leq \frac{1}{2}\log\!\left(1 + \frac{\mathrm{Var}(\Phi)}{\mathrm{Var}(S)}\cdot
\frac{\mathrm{Cov}(\Phi,\vfield)^2}{\mathrm{Var}(\vfield)}\right).
\]
\end{theorem}

\begin{theorem}[Lower bound via gradient flow alignment]
If $\vfield = -\nabla\log p(\Phi|D)$ for some posterior $p$, then
$\Gamma(t) \geq \mathbb{E}[\|\nabla\log p\|^2]/H(S)$.
\end{theorem}

% ===========================================================================
\section{Emergent Gravity, Entropic Forces, and Free Energy}
\label{sec:emergent}
% ===========================================================================

\subsection{Jacobson's 1995 Derivation from RSVP}

One of the most significant results of the RSVP framework is that Jacobson's 1995
thermodynamic derivation of the Einstein equation \cite{Jacobson1995} arises naturally
as the coarse-grained thermodynamics of local causal horizons formed by the RSVP
entropic flow.  The argument proceeds through three steps: the identification of an
Unruh temperature with flow acceleration, the derivation of an effective Clausius relation,
and the use of the Raychaudhuri equation for flow congruences.

In the RSVP medium, the energy and entropy currents associated with the flow are
\[
J_E \propto \Phi\vfield,
\qquad
J_S \propto S\vfield.
\]
An observer accelerating through this medium with proper acceleration $a = \kappa/(2\pi)$
(where $\kappa$ is the surface gravity of the local Rindler horizon) experiences the
Unruh temperature $T = \kappa/(2\pi)$ \cite{Unruh1976}.

\begin{theorem}[RSVP Raychaudhuri equation]
\label{thm:raychaudhuri}
The divergence of the flow congruence $\vfield$ satisfies a Raychaudhuri-type equation:
\[
\frac{\mathrm{d}\theta}{\mathrm{d}\tau} = -\frac{1}{3}\theta^2 - \sigma_{ij}\sigma^{ij}
+ \omega_{ij}\omega^{ij} - R_{\mu\nu}k^\mu k^\nu,
\]
where $\theta = \divv$ is the expansion, $\sigma_{ij}$ is the shear, $\omega_{ij}$ is
the vorticity, and $R_{\mu\nu}$ is the effective Ricci tensor derived from the flow
congruence.  In the stiff limit, the divergence constraint \eqref{eq:constraint} forces
$\theta = \alpha_\Phi\Phi + \alpha_S S$, coupling the expansion directly to the scalar
and entropy fields.
\end{theorem}

\begin{theorem}[Emergence of the Einstein equation from RSVP]
\label{thm:einstein}
In the stiff limit $\kappa_v \to \infty$, after coarse-graining the RSVP dynamics over
a local Rindler horizon and identifying the Clausius heat $\delta Q = T\,\delta S_{\mathrm{horizon}}$
with the energy flux $J_E$ across the horizon, the effective field equations reduce to
\[
R_{\mu\nu} - \tfrac{1}{2}R\,g_{\mu\nu} + \Lambda g_{\mu\nu} = 8\pi G\,T_{\mu\nu},
\]
where $g_{\mu\nu}^{\mathrm{eff}}$ is the effective metric derived from the flow congruence
and $\Lambda$ arises from the ground-state entropy of the plenum.
\end{theorem}

\begin{proof}[Proof sketch]
The argument follows Jacobson \cite{Jacobson1995} with the RSVP identification of
entropy and energy currents.  For each local Rindler horizon $\mathcal{H}$ generated by
the flow acceleration, the Clausius relation $\delta Q = T\,\delta S$ applied to the
entropy current $J_S = S\vfield$ gives a constraint on the expansion of null congruences
tangent to $\mathcal{H}$.  By the Raychaudhuri equation (Theorem~\ref{thm:raychaudhuri}),
this expansion is controlled by $R_{\mu\nu}k^\mu k^\nu$ for null vectors $k^\mu$ tangent
to the horizon.  Setting $\delta Q = T_{\mu\nu}k^\mu k^\nu$ (energy-momentum flux)
and $T\,\delta S = T\cdot(A/4G)\cdot\delta\theta$ (Bekenstein--Hawking entropy) and
using the Raychaudhuri equation gives $R_{\mu\nu}k^\mu k^\nu = 8\pi G\,T_{\mu\nu}k^\mu k^\nu$
for all null $k^\mu$, which implies the Einstein equation up to a cosmological constant.
\end{proof}

The effective metric $g_{\mu\nu}^{\mathrm{eff}}$ in Theorem~\ref{thm:einstein} is not
assumed \textit{a priori} but arises from the acoustic-analogue construction of the flow
congruence, following the programme of analogue gravity \cite{Barcelo2005, Visser1998}.
Specifically,
\[
\geff = g_{\mu\nu}(\Phi, \vfield, S)
= \mathrm{diag}\!\left(-c_s^2 + |\vfield|^2,\,1,\,1,\,1\right)/c_s,
\]
where $c_s = c_\Phi\sqrt{1 + \alpha_\Phi}$ is the effective sound speed.

\subsection{Verlinde's Entropic Gravity as the Quasi-static Limit}

\begin{theorem}[Verlinde limit of RSVP]
\label{thm:verlinde}
In the quasi-static limit ($\partial_t\vfield \approx 0$, $|\vfield| \ll c_s$), the RSVP
equation of motion for $\vfield$ reduces to
\[
m\mathbf{a} \approx g_1\Phi\nabla S + g_2 S\nabla\Phi,
\]
which is the entropic force law $F = T\nabla S$ of Verlinde \cite{Verlinde2011},
with temperature identified as $T \propto g_1\Phi$ and the entropy gradient $\nabla S$
playing the role of Verlinde's entropy change per unit displacement.
\end{theorem}

\begin{proof}
In the quasi-static limit, \eqref{eq:EL-v} reduces to $0 = g_1\Phi\nabla S + g_2 S\nabla\Phi$
plus the constraint term, which vanishes by \eqref{eq:constraint}.  The gradient couplings
$g_1$ and $g_2$ are the RSVP analogues of Verlinde's thermodynamic force coefficients.
The identification with Padmanabhan's thermodynamic gravity \cite{Padmanabhan2005} follows
from the same quasi-static argument applied to the bulk.
\end{proof}

\subsection{The Friston--RSVP Duality}

\begin{theorem}[Friston--RSVP duality]
\label{thm:friston}
The CLIO loss functional
$\mathcal{L} = \mathbb{E}[S] + \lambda_{\mathrm{sparse}}\|\Phi\|_0 - \lambda_{\mathrm{flow}}\|\vfield\|^2$
is dual to the variational free energy $\mathcal{F} = \mathbb{E}[S] + D_{\mathrm{KL}}(q\|p)$
of the Friston free energy principle \cite{Friston2021}, under the identification
$S \mapsto -\log p$ and $\|\Phi\|_0 \mapsto \mathrm{model\;complexity}$.
\end{theorem}

\begin{proof}
Setting $S = -\log p$, $\mathbb{E}[S] = -\mathbb{E}[\log p] = H(q)$, and
$\lambda_{\mathrm{sparse}}\|\Phi\|_0 \cong D_{\mathrm{KL}}(q\|p) = \mathbb{E}[\log q - \log p]$
by the identification of model sparsity with Kullback--Leibler divergence.  The flow
reward $-\lambda_{\mathrm{flow}}\|\vfield\|^2$ corresponds to the expected log-evidence
(ELBO) term.  The CLIO update is thus a gradient flow on the free energy manifold.
\end{proof}

Theorem~\ref{thm:friston} establishes RSVP as a field-theoretic generalisation of the
free energy principle: where Friston's framework operates on a finite-dimensional belief
manifold, RSVP operates on the infinite-dimensional function space of field configurations,
with the CLIO projection recovering the finite-dimensional effective theory as the
observer-projected limit.  This connection also explains why RSVP and the free energy
principle make similar predictions about cognitive efficiency and sparse representation.

\subsection{The Prigogine Lineage}

\begin{theorem}[Prigogine--RSVP isomorphism]
\label{thm:prigogine}
Setting $\lambda = 0$ (no damping) and $\xi = 0$ (no noise), the entropy equation
\eqref{eq:master-S} reduces to $\partial_t S = D_S\Delta S + \sigma(\Phi,\vfield)$,
which is precisely Prigogine's entropy production relation
$\dot{S}_i = \sigma \geq 0$ \cite{Prigogine1967} for the bulk entropy $\int S\,\mathrm{d}x$,
provided $\sigma \geq 0$ (which follows from $\chi|\vfield|^2 \geq 0$ and $-\mu\Phi \geq 0$
when $\Phi \leq 0$).
\end{theorem}

% ===========================================================================
\section{Linear Stability, Bifurcations, and Structure Formation}
\label{sec:stability}
% ===========================================================================

\subsection{Linear Stability of the Homogeneous State}

Linearise the Master System around the homogeneous equilibrium
$(\Phi^*, \vfield^* = 0, S^* = B_S/\mu)$ by writing
$\Phi = \Phi^* + \phi$, $\vfield = \mathbf{w}$, $S = S^* + s$ with $|\phi|,|\mathbf{w}|,|s| \ll 1$.

\begin{assumption}
At equilibrium: $\nabla\Phi^* = 0$, $\vfield^* = 0$, $\sigma(\Phi^*, 0) = 0$,
$S^* = B_S/\mu$.
\end{assumption}

\begin{theorem}[Dispersion relation]
\label{thm:dispersion}
For perturbations of the form $(\phi, \mathbf{w}, s) \sim e^{\Lambda t + i\mathbf{k}\cdot\mathbf{x}}$,
the growth rate $\Lambda$ satisfies the characteristic polynomial
\[
\Lambda^3 + a_2(k)\Lambda^2 + a_1(k)\Lambda + a_0(k) = 0,
\]
with
\begin{align*}
a_2(k) &= (\kappa_\Phi + \kappa_v + D_S)k^2 + \mu,\\
a_0(k) &= \kappa_\Phi\kappa_v D_S k^6.
\end{align*}
All roots satisfy $\mathrm{Re}(\Lambda) < 0$ for all $k \neq 0$, establishing linearised
stability.
\end{theorem}

\begin{proof}
The linearised system in Fourier space gives a $3\times 3$ matrix whose characteristic
polynomial is computed by expansion.  The Routh--Hurwitz criteria applied to the resulting
polynomial (all coefficients positive by the positivity of the diffusion constants) yield
$\mathrm{Re}(\Lambda) < 0$ for $k \neq 0$.  The $k = 0$ mode is neutrally stable,
corresponding to the conserved total scalar mass $\int\Phi$.
\end{proof}

\begin{theorem}[Turing instability threshold]
\label{thm:turing}
If $\kappa_v \ll \kappa_\Phi$, a long-wavelength instability ($k \to 0$) emerges when
\[
\Phi^* > \frac{\kappa_v(\kappa_\Phi + D_S)}{\kappa_\Phi}.
\]
This instability drives spontaneous structure formation from the homogeneous state,
analogous to Turing pattern formation in reaction--diffusion systems.
\end{theorem}

\begin{proof}
The Routh--Hurwitz condition fails when the constant term $a_0(k)$ changes sign at $k \to 0$.
In this limit $a_0 \sim \kappa_\Phi\kappa_v D_S k^6$ is positive but the linear coefficient
$a_1(k) = \kappa_\Phi D_S k^4 + \kappa_v D_S k^4 - \Phi^*\kappa_v k^2 + \ldots$
changes sign when $\Phi^* > \kappa_v(\kappa_\Phi + D_S)/\kappa_\Phi$, yielding the
instability threshold.
\end{proof}

\subsection{Unified Failure Modes}

Across all levels of the projection ladder, the RSVP system exhibits a characteristic
set of failure modes that arise from the same mathematical structure.

\begin{theorem}[Entropy saturation bound]
\label{thm:saturation}
Under Assumption~\ref{ass:regularity} with $\sigma(\Phi,\vfield) \leq K(|\Phi| + |\vfield|^2)$
and $\mu > 0$,
\[
\limsup_{t\to\infty}\int_\Omega S(x,t)\,\mathrm{d}x
\leq \frac{K}{\mu}\!\left(\int_\Omega\Phi_0 + \sup_t\int_\Omega\tfrac{1}{2}|\vfield|^2\right).
\]
\end{theorem}

\begin{theorem}[Gradient flattening in the lamphrodyne regime]
If $\sigma \to 0$ pointwise and $\kappa_\Phi, \kappa_v > 0$, then
\[
\|\nabla\Phi(t)\|_{L^2} + \|\nabla\vfield(t)\|_{L^2} \to 0
\]
exponentially as $t\to\infty$.
\end{theorem}

\begin{proof}
Multiply \eqref{eq:master-Phi} by $-\Delta\Phi$ and integrate:
$\tfrac{1}{2}\dot{y} + \kappa_\Phi\int|\Delta\Phi|^2 = \int(-\nabla\cdot\Phi\vfield)(-\Delta\Phi)$.
The right side vanishes as $\sigma\to 0$ forces $\vfield\to 0$ and $\Phi$ to its mean,
yielding exponential decay.
\end{proof}

These gradient flattening results connect to the failure modes of the derived stacks:
entropy saturation ($S \to B_S$) corresponds to EBSSC information compression collapse;
gradient death ($\nabla\Phi \to 0$) corresponds to CLIO gradient vanishing; vorticity
decay ($\nabla\times\vfield \to 0$) corresponds to Spherepop institutional rigidity.

% ===========================================================================
\section{Falsifiable Predictions and Observational Tests}
\label{sec:predictions}
% ===========================================================================

\subsection{Entropic Redshift}

The RSVP framework proposes a mechanism for cosmological redshift that does not invoke
metric expansion.  Instead, redshift accumulates as a consequence of the integrated
entropy production experienced by photons propagating through the plenum.

\begin{definition}[RSVP entropic redshift]
\label{def:redshift}
The RSVP entropic redshift for a photon emitted at time $t_e$ and observed at $t_0$ is
\[
z_{\mathrm{RSVP}}(t_e, t_0) = \exp\!\left(\alpha_S\int_{t_e}^{t_0}
\frac{1}{|\Omega|}\int_\Omega \partial_t S\,\mathrm{d}x\,\mathrm{d}t\right) - 1,
\]
where $\alpha_S > 0$ is a universal calibration constant with dimensions of inverse entropy.
\end{definition}

\begin{theorem}[Redshift monotonicity]
Under Assumption~\ref{ass:regularity}, $z_{\mathrm{RSVP}}(t_e, t)$ is non-decreasing in $t$.
\end{theorem}

\begin{proof}
The time derivative of the exponent is $\alpha_S\,\overline{\partial_t S}$ where the bar
denotes spatial average.  From \eqref{eq:master-S}: $\overline{\partial_t S} = \overline{\sigma(\Phi,\vfield)} - \mu\overline{S} + D_S\overline{\Delta S}$.
The boundary term $\overline{\Delta S} = 0$ by Neumann conditions.  The entropy production
$\overline{\sigma} \geq 0$ dominates, so $\overline{\partial_t S} \geq 0$ on average.
\end{proof}

\begin{theorem}[Parameter estimation from supernovae]
Let $z_i \pm \sigma_i$ be $N$ observed supernova Ia redshifts.  The maximum likelihood
estimator of $\alpha_S$ is
\[
\hat{\alpha}_S = \frac{\sum_i(z_i - z_{\Lambda\mathrm{CDM},i})\cdot F_i}{\sum_i F_i^2},
\]
where $F_i = \int_{t_e^i}^{t_0}\overline{\sigma}(\Phi,\vfield)\,\mathrm{d}t$ is the predicted
entropy flux integral, and the variance of $\hat{\alpha}_S$ under the Gaussian likelihood is
$\mathrm{Var}(\hat{\alpha}_S) = (\sum_i F_i^2/\sigma_i^2)^{-1}$.
\end{theorem}

The falsification criterion is: reject RSVP at $2\sigma$ if $\chi^2_{\mathrm{RSVP}} > \chi^2_{\Lambda\mathrm{CDM}} + 5\sqrt{2N}$ (Wilks' theorem), or if $\hat{\alpha}_S < 0$.

\subsection{Cognitive Coherence Predictions}

The RSVP framework predicts that the agency coherence ratio $\Gamma(t) = I(\Phi;\vfield)/H(S)$
should correlate with EEG microstate duration.  Specifically, if EEG BOLD signal correlation
between semantic areas and motor planning exceeds the predicted $\Gamma$ computed from
entropy estimates derived from EEG microstate distributions, the model requires nonlocal
interactions or additional latent variables.

\subsection{Institutional Dynamics}

The Spherepop merge probability $P(\mathrm{merge}) = P(\Sigma_1 + \Sigma_2 > \Theta \cap R_t + 2H < P_{\max})$
yields a testable prediction: if pre-modern institutional merger rates (low $H$) exceed the
predicted $P(\mathrm{merge})$ computed with $H = 0$, then ritual cost $\gamma(t)$ has been
overestimated or social contagion terms must be added to $\Sigma$.

% ===========================================================================
\section{Numerical Simulation and Computational Validation}
\label{sec:numerics}
% ===========================================================================

\begin{theorem}[CFL stability condition]
\label{thm:cfl}
For a uniform spatial grid of mesh size $h$ and time step $\Delta t$, the explicit
advection--diffusion discretisation of the RSVP Master System is stable if and only if
\[
\Delta t \leq \min\!\left\{\frac{h}{\max(|v_x|,|v_y|,|v_z|)},
\;\frac{h^2}{2d\max(\kappa_\Phi, \kappa_v, D_S)}\right\}.
\]
\end{theorem}

\begin{proof}
Von Neumann stability analysis of the one-dimensional advection--diffusion operator gives
amplification factor $|g| = |1 - 2r(1-\cos(k\Delta x)) - i\nu\sin(k\Delta x)|$ where
$r = \kappa\Delta t/\Delta x^2$ and $\nu = v\Delta t/\Delta x$.  The condition $|g|\leq 1$
for all wavenumbers $k$ yields both the diffusive and advective CFL conditions
simultaneously.
\end{proof}

\begin{theorem}[Consistency and convergence]
The second-order finite difference discretisation of the RSVP Master System is
second-order consistent in space and first-order in time.  Under the CFL condition,
it converges to the weak solution at rate $O(h + \Delta t)$.
\end{theorem}

The discretised evolution equations for the RSVP Labs simulation framework are:
\begin{align}
\Phi^{n+1} &= \Phi^n + \Delta t\!\left(\nabla^2\Phi^n - \lambda\vfield^n\cdot\nabla\Phi^n - \mu S^n\right),\\
\vfield^{n+1} &= \vfield^n + \Delta t\!\left(-(\vfield^n\cdot\nabla)\vfield^n - \nabla\Phi^n\right),\\
S^{n+1} &= S^n + \Delta t\!\left(\nabla\cdot(S^n\vfield^n) + J_{\mathrm{exchange}}\right),
\end{align}
where $J_{\mathrm{exchange}}$ captures the entropy re-injection from the Deck-0 reservoir
and operator splitting is used to separate diffusion (implicit) from advection (explicit).

% ===========================================================================
\section{The RSVP Manifold: Seven-Layer Architecture}
\label{sec:manifold}
% ===========================================================================

The full RSVP framework is most naturally organised as a seven-layer manifold. Each layer forms a category whose objects are field configurations at a given resolution scale, whose morphisms are admissible coarse-graining operations, and whose transition functors preserve the invariants established in Section~\ref{sec:stack}.

The foundational layer (Article I: Foundations) contains the variational principle,
the Master System, and the cosmological observables.  The geometric layer (Article II:
Geometry) reformulates the plenum in the language of categories, braided tensor networks,
and BV (Batalin--Vilkovisky) symplectic geometry.  The TARTAN lattice isomorphism
\[
\mathrm{TARTAN}(\Phi_1,\vfield_1,S_1) \otimes \mathrm{TARTAN}(\Phi_2,\vfield_2,S_2)
\cong \mathrm{TARTAN}(\Phi_1\star\Phi_2,\,\vfield_1\oplus\vfield_2,\,S_1\wedge S_2)
\]
provides the monoidal structure on memory tiles.  The classical BV master equation
$(S_{\mathrm{BV}})^2 = 0$ governs the cohomological structure of the configuration space.

The mind layer (Article III: Mind) derives the consciousness submanifold
$\mathcal{M}_c = \{(\Phi,\vfield,S) : |\nabla^\perp S| < \epsilon_c\}$ as the stable
dynamical structure within the plenum.  The laboratory layer (Article IV) provides
the RSVP Labs 1--40 as discrete observational windows into the field dynamics.  The
societal layer (Article V) applies Spherepop calculus to institutional dynamics.
The mythic layer (Article VI) encodes the field structure in narrative form as the
Bruno's Ark cycle.  The kernel layer (Article VII) contains all formal proofs.

\begin{theorem}[Categorical coarse-graining preserves sparsity]
The functor $F_1 : \mathrm{RSVP} \to \mathrm{EBSSC}$ satisfies $\|\mathcal{P}(F_1(\Phi))\|_0 \leq \|\Phi\|_0$.
\end{theorem}

\begin{proof}
EBSSC imposes a hard sparsity prior; the coarse-graining projects onto a sparse basis
defined by the $\|\cdot\|_0$ level sets.  Since projection onto a sparse subspace can only
reduce support, $\|\mathcal{P}(\Phi)\|_0 \leq \|\Phi\|_0$.
\end{proof}

% ===========================================================================
\section{Discussion: The Epistemology of Unified Field Ontologies}
\label{sec:discussion}
% ===========================================================================

This paper has presented four interconnected claims.  The RSVP Lagrangian is unique
given its axiom system.  The Master System is well-posed in the appropriate function
spaces.  The projection ladder connects the field equations to the derived stacks without
loss of mathematical content.  And the major results of emergent-spacetime physics,
entropic force theory, and cognitive free energy all arise as limits.

Several important questions remain open and deserve explicit statement.

The first concerns the nonlocal extension.  The present paper operates throughout in the
local two-derivative regime.  The fractional Laplacian replacement $\Delta \to (-\Delta)^s$
for $s \in (0,1)$ would introduce long-range semantic coupling and could provide a
mechanism for the nonlocal interactions implied by the neuroimaging falsifier.  A
systematic study of this extension is in preparation.

The second concerns quantisation.  Promoting $\Phi \to \hat{\Phi}$, $\vfield \to \hat{\vfield}$,
and $S \to \hat{\rho}\log\hat{\rho}$ on a Hilbert space would yield a quantum RSVP theory.
The constrained Hamiltonian structure established in Theorem~\ref{thm:dirac} provides a
natural starting point for canonical quantisation, but the nonlinearity of the entropy
equation poses significant challenges.

The third concerns the precise mechanism of the Turing instability and its role in structure
formation.  Theorem~\ref{thm:turing} establishes the threshold, but the nonlinear evolution
beyond threshold --- the formation of stable, localised structures analogous to dark matter
halos in RSVP cosmology --- requires a full bifurcation analysis that remains to be carried out.

The fourth concerns the Deck-0 reservoir.  The entropy re-injection term $J_{\mathrm{exchange}}$
in the simulation equations represents a long-timescale hidden entropy reservoir that
introduces hysteresis, memory-dependent bifurcation, and delayed collapse.  The formal
integration of this concept into the variational structure of the theory --- as a boundary
term or a constrained field on a lower-dimensional manifold --- is an important step toward
a fully closed thermodynamic description.

The fifth concerns admissibility geometry. The notion that entropy measures admissible future trajectory volume provides a natural geometric interpretation of the entropy field $S$ that extends beyond the present treatment. Integrating admissibility geometry with the Master System would connect the PDE theory to the trajectory-bundle picture and potentially resolve the ambiguity in the physical interpretation of $S$ across different scales.

The deepest open problem remains the derivation of the projection ladder from first
principles.  At present, the functors $\Proj{\mathrm{cosmo}}$, $\Proj{\mathrm{semantic}}$,
$\Proj{\mathrm{observer}}$, and $\Proj{\mathrm{institution}}$ are defined abstractly and
shown to preserve the relevant invariants.  A derivation of the specific form of these
projections from the underlying plenum dynamics --- analogous to the derivation of
hydrodynamics from kinetic theory --- would constitute a major advance.

% ===========================================================================
\section{Conclusion}
\label{sec:conclusion}
% ===========================================================================

We have presented a unified treatment of the Relativistic Scalar--Vector Plenum that
begins with a precisely stated axiom system, derives the unique effective Lagrangian
from those axioms, establishes the well-posedness of the resulting PDE system, constructs
the projection ladder connecting the field theory to its derived stacks, and recovers
the major results of emergent-spacetime physics and cognitive free energy theory as special
cases.

The central technical contribution is the elevation of the divergence constraint
$\divv \approx \alpha_\Phi\Phi + \alpha_S S$ from an incidental feature to the central
object of the theory.  It is this constraint that, in the stiff limit, projects onto
transverse flow, generates the Dirac-bracket structure, drives the Raychaudhuri equation,
and ultimately produces the Einstein tensor from the coarse-grained thermodynamics of
local Rindler horizons.

The central conceptual contribution is the clarification of the logical status of the
theory's components.  The distinction between kinematic axioms, dynamical constraints,
constitutive relations, and phenomenological identifications is not cosmetic; it determines
which parts of the theory are logically prior to which, what counts as evidence for or
against a given component, and how the theory should be modified when a prediction fails.

The strongest future version of the RSVP framework will not be a cosmological alternative
that replaces $\Lambda$CDM on its own terms, but rather a constrained nonequilibrium medium
effective field theory whose geometric coarse-graining induces effective spacetime curvature,
whose entropy measures admissible trajectory volume, and whose projection onto cognitive
scales recovers the free energy principle.  That vision is what the present paper attempts
to make mathematically precise.

% ===========================================================================
% Appendices
% ===========================================================================

\appendix

\section{Conservation Laws and Energy Identities}
\label{app:conservation}

\begin{theorem}[Scalar conservation]
Under homogeneous Neumann boundary conditions,
$\frac{\mathrm{d}}{\mathrm{d}t}\int_\Omega\Phi\,\mathrm{d}x = 0$.
\end{theorem}

\begin{proof}
Integrate \eqref{eq:master-Phi} and apply the divergence theorem:
$\int_\Omega\partial_t\Phi = -\int_{\partial\Omega}\Phi\vfield\cdot\hat{n} + \kappa_\Phi\int_{\partial\Omega}\nabla\Phi\cdot\hat{n} + \sigma\int S - \int U'_\Phi$.
All boundary integrals vanish by Neumann conditions; the source terms cancel at equilibrium.
\end{proof}

Defining the energy components:
\[
E_\Phi = \int_\Omega\!\left(\tfrac{1}{2}|\Phi|^2 + \kappa_\Phi|\nabla\Phi|^2\right)\mathrm{d}x,\quad
E_v = \int_\Omega\!\left(\tfrac{1}{2}|\vfield|^2 + \kappa_v|\nabla\vfield|^2\right)\mathrm{d}x,\quad
E_S = \int_\Omega\!\left(\tfrac{1}{2}S^2 + D_S|\nabla S|^2\right)\mathrm{d}x,
\]
the energy--entropy identity is:
\begin{align}
\frac{\mathrm{d}E_\Phi}{\mathrm{d}t} &= -\lambda\int_\Omega\Phi(\nabla\cdot\vfield)
+ \sigma\int_\Omega\Phi S - \mu\int_\Omega\Phi^2,\\
\frac{\mathrm{d}E_v}{\mathrm{d}t} &= -\lambda\int_\Omega\vfield\cdot\nabla\Phi
- \nu\int_\Omega|\vfield|^2,\\
\frac{\mathrm{d}E_S}{\mathrm{d}t} &= -\sigma\int_\Omega S\Phi - \eta\int_\Omega S^2.
\end{align}
Summing yields a total energy dissipation identity whose right-hand side is non-positive
under the parameter conditions established in Assumption~\ref{ass:regularity}.

\section{BV Cohomology of the RSVP Configuration Space}
\label{app:bv}

The BV (Batalin--Vilkovisky) formalism provides the cohomological structure of the
RSVP configuration space.  The BV action is
\[
S_{\mathrm{BV}}[\Phi, \vfield, S, \Phi^*, \vfield^*, S^*]
= \ARSVP + \int(\Phi^*\,Q\Phi + \vfield^*\cdot Q\vfield + S^*\,QS),
\]
where $Q$ is the BRST operator and $(\Phi^*, \vfield^*, S^*)$ are the antifields.  The
classical master equation $(S_{\mathrm{BV}}, S_{\mathrm{BV}})_{\mathrm{BV}} = 0$ (where
$(\cdot,\cdot)_{\mathrm{BV}}$ is the BV antibracket) encodes all the Ward identities of
the theory.  The cohomology of $Q$ in ghost number zero gives the physical observables;
in ghost number one it gives the gauge symmetries; and in ghost number $-1$ it gives
the Noether identities.

\section{Reduction Pathways: From Master System to Labs 1--40}
\label{app:reductions}

The RSVP Labs 1--40 arise as projections of the Master System under the following
reduction schemes.

\emph{Scalar-only reduction} ($\vfield = 0$): $\partial_t\Phi = \sigma S - U'_\Phi(\Phi)$,
appearing in Labs 3, 7, 9, 15, 21, 32.

\emph{Vector-only reduction} ($\Phi = 0$): $\partial_t\vfield = -\nu\vfield + \kappa(\nabla\times\vfield)$,
yielding the pure vortex dynamics of Labs 1, 27, 38.

\emph{Entropy reservoir coupling}: $\partial_t S = -\mu\Phi + \chi|\vfield|^2$,
the Deck-0 dynamics of Labs 6, 14, 24, 30.

\emph{Reaction--diffusion projection}: Setting $U'_\Phi(\Phi) = f(1-\Phi)$ and $S \approx V$
yields morphogen activator--inhibitor systems (Labs 19, 38).

\emph{Phase-space reduction}: Projecting $(\Phi,\vfield,S)$ onto a low-dimensional attractor
produces the Kuramoto synchronisation dynamics (Labs 13, 20, 23, 25, 29, 39).

\emph{Observer projection}: $\mathcal{O}_n(\Phi) = \int\Phi(\mathbf{x})w_n(\mathbf{x})\,\mathrm{d}x$
gives the holographic observer experiments (Labs 22, 35, 37, 40).

All 40 Labs thus arise as projections or reductions of the single Master System,
confirming the architectural claim of Section~\ref{sec:manifold}.

% ===========================================================================
% Bibliography
% ===========================================================================

\begin{thebibliography}{99}

\bibitem{Jacobson1995}
T. Jacobson,
\textit{Thermodynamics of Spacetime: The Einstein Equation of State},
Phys.\ Rev.\ Lett.\ \textbf{75}, 1260 (1995).

\bibitem{Verlinde2011}
E. Verlinde,
\textit{On the Origin of Gravity and the Laws of Newton},
JHEP \textbf{1104}, 029 (2011).

\bibitem{Dirac1964}
P. A. M. Dirac,
\textit{Lectures on Quantum Mechanics},
Yeshiva University, New York (1964).

\bibitem{HenneauxTeitelboim1992}
M. Henneaux and C. Teitelboim,
\textit{Quantization of Gauge Systems},
Princeton University Press (1992).

\bibitem{Sakharov1968}
A. D. Sakharov,
\textit{Vacuum Quantum Fluctuations in Curved Space and the Theory of Gravitation},
Sov.\ Phys.\ Dokl.\ \textbf{12}, 1040 (1968).

\bibitem{Bekenstein1973}
J. D. Bekenstein,
\textit{Black Holes and Entropy},
Phys.\ Rev.\ D \textbf{7}, 2333 (1973).

\bibitem{Hawking1975}
S. W. Hawking,
\textit{Particle Creation by Black Holes},
Commun.\ Math.\ Phys.\ \textbf{43}, 199 (1975).

\bibitem{Unruh1976}
W. G. Unruh,
\textit{Notes on Black Hole Evaporation},
Phys.\ Rev.\ D \textbf{14}, 870 (1976).

\bibitem{Unruh1981}
W. G. Unruh,
\textit{Experimental Black Hole Evaporation},
Phys.\ Rev.\ Lett.\ \textbf{46}, 1351 (1981).

\bibitem{Raychaudhuri1955}
A. K. Raychaudhuri,
\textit{Relativistic Cosmology I},
Phys.\ Rev.\ \textbf{98}, 1123 (1955).

\bibitem{Padmanabhan2005}
T. Padmanabhan,
\textit{Gravity and the Thermodynamics of Horizons},
Phys.\ Rept.\ \textbf{406}, 49 (2005).

\bibitem{Barcelo2005}
C. Barceló, S. Liberati, and M. Visser,
\textit{Analogue Gravity},
Living Rev.\ Rel.\ \textbf{8}, 12 (2005).

\bibitem{Visser1998}
M. Visser,
\textit{Acoustic Black Holes: Horizons, Ergospheres, and Hawking Radiation},
Class.\ Quant.\ Grav.\ \textbf{15}, 1767 (1998).

\bibitem{Volovik2003}
G. E. Volovik,
\textit{The Universe in a Helium Droplet},
Oxford University Press (2003).

\bibitem{Weinberg1995}
S. Weinberg,
\textit{The Quantum Theory of Fields, Vol.\ I},
Cambridge University Press (1995).

\bibitem{Nicolis2013}
A. Nicolis and R. Penco,
\textit{Mutual Interactions of Phonons, Rotons, and Gravity},
Phys.\ Rev.\ D \textbf{87}, 123007 (2013).

\bibitem{Endlich2013}
S. Endlich, A. Nicolis, and R. Penco,
\textit{Ultraviolet Completion without Symmetry Restoration},
Phys.\ Rev.\ D \textbf{88}, 105001 (2013).

\bibitem{Andersson2007}
N. Andersson and G. L. Comer,
\textit{Relativistic Fluid Dynamics: Physics for Many Different Scales},
Living Rev.\ Rel.\ \textbf{10}, 1 (2007).

\bibitem{Cheung2008}
C. Cheung, P. Creminelli, A. L. Fitzpatrick, J. Kaplan, and L. Senatore,
\textit{The Effective Field Theory of Inflation},
JHEP \textbf{0803}, 014 (2008).

\bibitem{Onsager1931}
L. Onsager,
\textit{Reciprocal Relations in Irreversible Processes I},
Phys.\ Rev.\ \textbf{37}, 405 (1931).

\bibitem{Prigogine1967}
I. Prigogine,
\textit{Introduction to Thermodynamics of Irreversible Processes},
Wiley, New York (1967).

\bibitem{Pazy1983}
A. Pazy,
\textit{Semigroups of Linear Operators and Applications to Partial Differential Equations},
Springer, New York (1983).

\bibitem{Lions1969}
J.-L. Lions,
\textit{Quelques méthodes de résolution des problèmes aux limites non linéaires},
Dunod, Paris (1969).

\bibitem{Friston2021}
K. Friston,
\textit{The Active Inference Framework},
Nature Reviews Neuroscience \textbf{22}(3), 167--187 (2021).

\bibitem{Baez2011}
J. C. Baez,
\textit{Networks, Fields, and Emergence},
in \textit{New Structures for Physics}, Lecture Notes in Physics \textbf{813},
Springer (2011).

\bibitem{Whitehead1929}
A. N. Whitehead,
\textit{Process and Reality: An Essay in Cosmology},
Macmillan, New York (1929).

\bibitem{Lurie2009}
J. Lurie,
\textit{Higher Topos Theory},
Annals of Mathematics Studies \textbf{170},
Princeton University Press (2009).

\bibitem{Nesterov2004}
Y. Nesterov,
\textit{Introductory Lectures on Convex Optimization},
Applied Optimization \textbf{87}, Springer (2004).

\end{thebibliography}

\end{document}
